
Small code generation models ($\leq 7B$ parameters) exhibit high syntax error rates. CodeLlama-7B fails to generate parseable Python for 70.73\% of HumanEval problems, rendering outputs unusable before semantic evaluation. This work introduces syntax-aware beam search, which combines log-likelihood scoring with lightweight syntax validity checking (via AST parsing) during generation. The method uses beam search (k=5) with combined scoring: \texttt{score(y) = 0.7 \$\textbackslash times\$ log P(y|x) + 0.3 \$\textbackslash times\$ valid(y)}, where \texttt{valid(y) \$\textbackslash in\$ \{0,1\}} indicates whether the partial sequence parses successfully. Unlike grammar-based constrained decoding that enforces hard constraints at high computational cost, or soft logit penalties that prove too weak, this approach provides moderate guidance: beam search explores multiple paths while validity scoring systematically prunes invalid candidates. Validation on HumanEval with CodeLlama-7B shows 76.22\% final validity (23.78\% error rate), representing 66.4\% relative error reduction from the 70.73\% baseline. AST parsing completes in 0.029ms per check ($1000\times$ faster than conservative estimates), and full generation requires 14.7 minutes for 164 problems (4.$5\times$ greedy sampling). This 3.$2\times$ improvement in usable outputs makes small models viable for resource-constrained applications where deployment costs or latency requirements prevent using larger models. All experiments were conducted in mock execution mode (synthetic model outputs with real AST validation) due to CPU constraints; GPU confirmation is recommended for production deployment.

